\scanpage{178}
% archive sheet number: 177
dire que $\underline{\unsure{\mathscr G|\mathcal G}}$ contient la représentation triviale dans $k$,
signifie que $1\in A_0$. \deleted{Dire que} Si \deleted{\unsure{donc}}
ces conditions sont \unsure{vérifiées|satisfaites}, \inserted{(on dit que $\underline{\mathscr G}$ est un \uline{\unsure{clan|classe} de représentations})} \deleted{\unsure{on a une}} on a une
$K$-catégorie, dont la \deleted{\unsure{coalg.}} \inserted{algèbre} des coéfficients est
$A_0$. Le \uline{groupe structural} est \inserted{donc} l'ensemble des
caractères de $A_0$, soit $G'$, et on a $G\unsure{-|\to}G'$.
\deleted{\unsure{Donc}} \inserted{Par définition} \uline{la dualité de Tannaka\unsure{.|,} vaut si et
seulement si $G\to G'(A_0)$ est bijectif},
i.e. \uline{$G$ est exactement le groupe des
caractères de $A_0$}, i.e. $A_0$ sépare les $\mathrm p^{\mathrm{ts}}$
de $G$, et tt caractère de $A_0$ provient de $G$.

A quelles condition la structure sur $G$ définie
par la donnée de $\underline{\mathscr G}$, i.e. de $A_0$, est
elle une structure de \uline{groupe algébrique
linéaire}$\unsure{??|?}$
\deleted{\unsure{i.e. si applicable à la situation}}
\deleted{\unsure{définie par}} \inserted{Il faut, et il suffit que}
$A_0$ \uline{soit à engendrement fini} sur $k$.

% slanted
\margin{de nouveau, il est inutile de se restreindre au cas d'un $A_0$ défini sur $G$, un groupe, mais on \unsure{peut} partir d'un $A_0$ quelconque.}
On sait que il revient au $=$ que
\uline{$R$ soit à ``engendrement fini''} en tant
qu'ensemble de représentations, i.e. qu'il
existe une $V\in R$ telle que \deleted{tout}
$W\in R$ soit contenue dans un $\overset{\unsure{n}}{\otimes}V$ ;
Nécessaire, car si les $f_i\in A_0$ sont des
générateurs de $A_0$, et les $V_i$ des $\in R$
